\documentclass[a4paper,11pt,french]{report} 
\usepackage{mathrsfs}
\usepackage{amsmath}% includes amsmath
\usepackage{amsthm}
\usepackage{hyperref}
\usepackage[tmargin=2cm,rmargin=1in,lmargin=1in,margin=0.85in,bmargin=2cm,footskip=.2in]{geometry}

\usepackage[default]{fontsetup}
\usepackage{fontawesome}
\renewcommand{\qedsymbol}{\faSmileO}
\theoremstyle{plain}
\newtheorem{theorem}{Théorème}[chapter]
\newtheorem{lemma}[theorem]{Lemme}
\newtheorem{proposition}[theorem]{Proposition}
\newtheorem{claim}[theorem]{Claim}
\newtheorem{corollary}[theorem]{Corollaire}

\theoremstyle{definition}
\newtheorem{definition}{Définition}[chapter]
\newtheorem*{notation}{Notation}
\newtheorem*{question}{Question}
\newtheorem*{answer}{Answer}
\newtheorem*{reminder}{Rappel}
\newtheorem{example}{Exemple}[chapter]
\newtheorem{remark}[theorem]{Remarque}
\newtheorem{problem}{Problème}[chapter]
\newtheorem{exercise}{Exercice}[chapter]
\title{Exos etoile}
\author{Nathalie Sarraute}
\begin{document}
\maketitle
\tableofcontents
\chapter*{References, references, references}
\begin{theorem}\label{thm22}
Soit $R$ un anneau alors il existe un anneau $R'\supseteq R$ tel que $R$ est un sous-anneau de $R'$ et $1_{R'}=1_{R}$ et ces conditions sont satisfaites:

Il existe $x \in R\setminus R$ tel que
\begin{enumerate} 
\item  $\forall a\in R:ax=xa$
\item Si pour $a_{0},\dots,a_{n}\in R:a_{0}+a_{1}x+\dots+a_{n}x^{n}=0$ alors tous les $a_{i}$ sont egaux a 0.
\end{enumerate}
On appele $x$ une variable ou indeterminee de $R$. 
\end{theorem}
\chapter{Serie 1}
\begin{exercise}[Exo plus]
Soit $R$ un anneau et $\alpha \in Z(R)$ un element du centre de $R$. Montrer que l'application

\begin{align}
\phi : R[x] & \to R \\
f(x) & \mapsto f(\alpha)
\end{align}

est un morphisme d'anneuax surjectif.

Soit $R=\mathbf{Z}^{2\times 2}$, trouver $\alpha \in R$ (qui n'appartient donc pas a $Z(R)$) et deux polynomes $g,h\in R[x]$ tel que
$$
g(\alpha)h(\alpha)\neq (gh)(\alpha)
$$
\end{exercise}
\begin{proof}[Solution]
Pour montrer que $\phi$ est un homomorphisme d'anneaux on verifie ces trois conditions $\forall f,g\in R[x]$
\begin{enumerate}
\item $\phi(1)=1$
\item $\phi(g+f)=\phi(g)+\phi(f)$
\item $\phi(fg)=\phi(f)\phi(g)$
\end{enumerate}
on verifie:
\begin{enumerate}
\item $\phi(1)=1$ car le polynome constant est envoye sur lui meme par $\phi$
\item \label{exo1} On prend $f,g\in R[x]$ et on ecrit $f(x)=\sum_{i=0}^{\deg f}a_{i}x^{i}$ et $g(x)=\sum_{j=0}^{\deg g}b_{j}x^{j}$ ou $a_{i},b_{j}\in R$. 
\begin{align}
\phi(f(x)+g(x)) & =\phi\left( \sum_{i=0}^{\deg f} a_{i}x^{i}+\sum_{j=0}^{\deg g} b_{j}x^{j}  \right) \\
 & =\phi\left( \sum_{k=0}^{\max\{ \deg f,\deg g \}} (a_{k}+b_{k})x^{k} \right) \\
 & =\sum_{k=0}^{\max\{ \deg f,\deg g \}}(a_{k}+b_{k})\alpha^{k} \\
 & =\sum_{i=0}^{\deg f} a_{i}\alpha^{i}+\sum_{j=0}^{\deg g} b_{j}\alpha^{j}  \\
 & =\phi(f(x))+\phi(g(x)) \\
\end{align}
\item On reprend $f,g$ comme avant et on calcule 
\begin{align}
\phi(f(x)g(x)) & =\phi\left( \sum_{i=0}^{\deg f} a_{i}x^{i} \sum_{j=0}^{\deg g} b_{j}x^{j} \right) \\
 & =\phi\left( \sum_{i=0}^{\deg f} \sum_{j=0}^{\deg g} a_{i}b_{j}x^{i+j} \right) \\
 & =\sum_{i=0}^{\deg f} \sum_{j=0}^{\deg g} a_{i}b_{j}\alpha^{i+j} \label{exo2}\\
 & =\left( \sum_{i=0}^{\deg f } a_{i}\alpha^{i} \right)\left( \sum_{j=0}^{\deg g } b_{j}\alpha^{j} \right) \\
 & =\phi(f(x))\phi(g(x))
\end{align}
\end{enumerate}
ou \ref{exo2} est obtenu car $\alpha \in Z(R)$. Donc $\phi$ est un morphisme d'anneaux qui est surjectif a cause de \ref{exo1}.

Soit $R=\mathbf{Z}^{2\times2}$, on prend $\alpha=\begin{pmatrix}1 & 1 \\ 0 & 1\end{pmatrix}$, $g(x)=\begin{pmatrix}1 & 1  \\ 0 & 1 \end{pmatrix}x$ et $h(x)=\begin{pmatrix} 1 & 0 \\ 1 & 1\end{pmatrix}x$ on a
$$
\begin{pmatrix}
1 & 1 \\
0 & 1
\end{pmatrix}\begin{pmatrix}
1 & 0 \\
1 & 1
\end{pmatrix}=\begin{pmatrix}
2 & 1 \\
1 & 1
\end{pmatrix}
$$
et 
$$
\begin{pmatrix}
1 & 0 \\
1 & 1
\end{pmatrix}\begin{pmatrix}
1 & 1 \\
0 & 1
\end{pmatrix}=\begin{pmatrix}
1 & 1 \\
1 & 2
\end{pmatrix}
$$
Donc ces matrices ne commutent pas. On obtient
$$
g(\alpha)=\begin{pmatrix}
1 & 2 \\
0 & 1
\end{pmatrix}\quad f(\alpha)=\begin{pmatrix}
1 & 1 \\
1 & 2
\end{pmatrix}
$$
donc
$$
g(\alpha)f(\alpha)=\begin{pmatrix}
3 & 5 \\
1 & 2
\end{pmatrix}
$$
mais
$$
g(x)h(x)=\begin{pmatrix}
2 & 1 \\
1 & 1 
\end{pmatrix}x^{2}
$$
que si on evalue en $\alpha$
$$
(gh)(\alpha)=\begin{pmatrix}
2 & 5 \\
1 & 3
\end{pmatrix}
$$
qui est different de $g(\alpha)f(\alpha)$
\end{proof}
\begin{exercise}[Exo etoile]
Soit $R$ un anneau.
\begin{enumerate}
\item Montrer que l'element $1$ est unique.
\item Montrer qu'un element inversible $r\in R^{*}$ n'est pas un diviseur de zero.
\item Montrer que deux polynomes 
$$
p(x)=a_{0}+a_{1}x+\dots+a_{n}x^{n}\quad\text{et}\quad q(x)=b_{0}+b_{1}x+\dots+b_{m}x^{m}\in R[x]
$$
sont egaux si et seulement si $a_{i}=b_{i}$ pour tous $i$.
\item Soit $R^{n\times n}$ l'anneau des matrices $n\times n$ sur $R$. Montrer que le centre de $R^{n\times n}$ est $Z(R^{n\times n})=\{ aI_{n}:a\in Z(R) \}$.
\end{enumerate}
\end{exercise}
\begin{proof}[Solution]
\begin{enumerate}
\item Supposons qu'il existe $1'\in R$ tel que $\forall a\in R$ on a : $1'a=a1'=a$. Prenons $a=1$ alors $1=11'=1'1=1'$ car $1$ est l'unite de $R$.
\item Prenons $r\in R^{*}$ et supposons qu'il existe $a\in R$ tel que $ra=0$ alors:
\begin{align*}
r^{-1}ra & =r^{-1}0 \\
a & =0
\end{align*}
donc $r$ ne peut pas etre un diviseur de zero.
\item Le sens $\Leftarrow$ est trivial car si deux polynomes ont les memes coefficients alors ce sont les memes polynomes.

Pour le sens $\Rightarrow$ on utilise \ref{thm22} qui nous dit que le seul polynome egal a 0 est le polynome dont tous les coefficients valent 0. Comme $p(x)-q(x)=0$ on a que $a_{i}-b_{i}=0$ pour tous $i$ donc $a_{i}=b_{i}$.
\item On procede par double inclusion: prenons $aI_{n}\in Z(R^{n\times n})$ alors $\forall A\in R^{n\times n}$ on a $aI_{n}A=aA$ et comme $a\in Z(R)$ on a l'equivalence $aA=Aa$ que montre que $aI_{n}A=A(aI_{n})$ pour montrer l'inclusion $\supseteq$.

Pour montrer l'inclusion $\subseteq$ on va utiliser les matrices elementaires $E^{ij}$. 
\begin{reminder}
\begin{itemize}
\item Pour $A\in R^{n\times n}$, la matrice $AE^{ij}$ est une matrice nulle sauf la $j$-eme colonne qui est egale a la $i$-eme colonne de $A$.
\item Pour $E^{ij}A$ la matrice est nulle sauf la $i$-eme ligne de qui est egale a la $j$-eme ligne de $A$.
\end{itemize}
\end{reminder}
Prenons $C\in Z(R^{n\times n})$. On considere les matrices $CE^{ii}$ et $E^{ii}C$ qui sont egales car $C\in Z(R^{n\times n})$ donc $C$ est une matrice diagonale qu'on note $C=\mathrm{diag}(c_{11},\dots,c_{nn})$. On compare les coefficients sur la diagonale qu'on denote par $(CE^{ij})_{ij}=c_{ii}$ et $(E^{ij}C)_{ij}=c_{jj}$ et comme ces matrices sont egales on a que $c_{ii}=c_{jj}$. On sait alors que $C=cI_{n}$ pour $c\in R$. Donc on a
\begin{align*}
(ac)I_{n} & =(aI_{n})(cI_{n}) \\
 & =(cI_{n})(aI_{n}) \\
 & =(ca)I_{n} 
\end{align*}

donc on a montre que $c\in Z(R)$.
\end{enumerate}
\end{proof}

\chapter{Serie 2}
\begin{exercise}[Exo plus]
Soit $K$ un corps et $f(x)\in K[x]$ un polynome de degre 3. Montrer que $f(x)$ est irreductible si et seulement si $f(x)$ n'a pas de racines dans $K$.
\end{exercise}
\begin{proof}[Solution]
\begin{definition}
Un polynome $g(x)\in K[x]$ est irreductible si
\begin{enumerate}
\item $\deg g\ge 1$
\item Si $g(x)=p(x)q(x)$ avec $p(x),q(x)\in K[x]$ alors $\deg p=0$ ou $\deg q=0$.
\end{enumerate}
\end{definition}
On montre le sens $\Rightarrow$ par contraposition. Supposons que $f(x)$ admet une racine $\alpha \in K$ (c-a-d $f(\alpha)=0$), alors $(x-\alpha)\mid f(x)$ qui implique que $\exists q(x)\in K[x]$ tel que $(x-\alpha)q(x)=f(x)$ donc $f(x)$ n'est pas irreductible.

Pour montrer le sens $\Leftarrow$ on suppose que $f(x)$ est reductible et de degre 3 alors il existe $\beta,\gamma \in K$ avec $\beta\neq 0$ tel que $(\beta x+\gamma)\mid f(x)$ alors $-\frac{\gamma}{\beta}$ est une racine de $f(x)$ qui pose une contradiction.
\end{proof}
\begin{exercise}[Exo etoile]
Soit $K$ un corps et $A\in K^{n\times n}$ une matrice inversible. Posons $\tilde{A}:=(\mathrm{com}A)^{\mathsf{T}}$ la transposee de la comatrice de $A$.
\begin{enumerate}
\item Montrer que $\det(A)\neq 0$.
\item Montrer que $A^{-1}=\det(A)^{-1}\tilde{A}$.
\item Montrer que pour toute matrice $A\in \mathbf{Z}^{n\times n}$, il existe $B\in \mathbf{Z}^{n\times n}$ telle que $AB=I$ si et seulement si $\det(A)=\pm1$.
\end{enumerate}
\begin{reminder}
$\mathrm{com}(A)_{ij}=(-1)^{i+j}\det(A_{ij})$ ou $A_{ij}$ est la matrice dont on a enlevee la $i$-eme ligne et $j$-eme colonne.
\end{reminder}
\end{exercise}
\begin{proof}[Solution]
\begin{enumerate}
\item Comme $A$ est inversible on a que $AA^{-1}=I_{n}$ alors $\det(AA^{-1})=1$ donc par proprietes du determinant on a $\det(A)\det (A^{-1})=1$ et $\det(A)\neq 0$.
\item Pour conclure, il suffit de montrer que $A\tilde{A}=\det(A)I_{n}$ (par unicite de l'inverse d'une matrice)( Pour etre precis, il faut aussi demontrer que $\tilde{A}A=\det(A)I_{n}$). Considerons les coefficients diagonaux de $A\tilde{A}$, pour tout $i\in \{ 1,\dots ,n \}$ on a 

\begin{align*}
(A\tilde{A})_{ii} & =\sum_{j=1}^{n} a_{ij}\tilde{a}_{ji} \\
 & =\sum_{j=1}^{n} a_{ij}(\mathrm{com}A)_{ij} \\
 & =\sum_{j=1}^{n} a_{ij}(-1)^{i+j}\det(A_{ij}) \\
 & =\det A.
\end{align*}

Considerons les coefficients hors diagonal. Pour $i\neq j\in \{ 1,\dots,n \}$ on a

\begin{align*}
(A\tilde{A})_{ij} & =\sum_{k=1}^{n} a_{ik}\tilde{a}_{kj} \\
 & =\sum_{k=1}^{n} a_{ik}(\mathrm{com}A)_{jk} \\
 & =\sum_{k=1}^{n} a_{ik}(-1)^{j+k}\det(A_{jk}) \\
 & =\text{le determinant de la matrice }A\text{ dont on remplace la }\\&j\text{-eme ligne par une copie de la }i\text{-eme ligne} \\
 & =0
\end{align*}
car $\det$ est une forme alternee.
\item Pour montrer le sens $\Rightarrow$, on suppose qu'il existe $B\in \mathbf{Z}^{n\times n}$ tel que $AB=I_{n}$ donc $\det(AB)=1$ et comme dans $\mathbf{Z}$ les seuls diviseurs de 1 sont 1 et $-1$ alors on observe que $\det(A)=\pm1$.

On montre l'autre sens $\Leftarrow$, on suppose que $\det(A)=\pm1$. Par les memes techniques que dans 2 on voit que $A\tilde{A}=\det(A)I_{n}$ avec $\tilde{A}=(\mathrm{com}A)^{\mathsf{T}}\in \mathbf{Z}^{n\times n}$. Comme $\det(A)=\pm 1$ on conclut que $B=\det(A)^{-1}\tilde{A}\in \mathbf{Z}^{n\times n}$.
\end{enumerate}
\end{proof}

\chapter{Serie 3}

\begin{exercise}[Exo plus]
Soit $V$ un espace vectoriel reel et soit $B=\{ v_{1},\dots,v_{4} \}$ une base de $V$.
\begin{enumerate}
\item Soit $f$ l'endomorphsime defini par
$$
f(v_{1})=v_{1}-v_{2},~f(v_{2})=2v_{2}-6v_{3},~f(v_{3})=-2v_{1}+2v_{2},~f(v_{4})=v_{2}-3v_{3}+v_{4}
$$
Ecrivez la matrice $A_{B}$ de l'application $f$ dans la base $B$. Est-ce que $f$ est inversible? Si oui, ecrivez la matrice $A^{-1}_{B}$ de l'application inverse $f^{-1}:V\to V$.
\end{enumerate}
\end{exercise}
\begin{proof}[Solution]
On construit la matrice $A_{B}$:
$$
A_{B}=\begin{pmatrix}
1 & 0 & -2 & 0 \\
-1 & 2 & 2 & 1 \\
0 & -6 & 0 & -3 \\
0 & 0 & 0 & 1
\end{pmatrix}
$$
on peut observer que il y a une relation de dependence entre la troisieme et premiere colonne donc $A_{B}$ n'est pas inversible.
\end{proof}

\begin{exercise}[Exo etoile]
Soit $K$ un corps et $f,g\in K[x]$ deux polynomes pas tous les deux nuls. Considerons l'ensemble des diviseurs communs a $f$ et $g$:
$$
D_{f,g}=\{ d\in K[x]:d\mid f , d\mid g\}
$$
\begin{enumerate}
\item Montrer qu'il existe un unique polynome $d\in D_{f,g}$ unitaire et de degre maximal.
\item Montrer que $d=\gcd(f,g)$.
\end{enumerate}
\end{exercise}
\begin{proof}[Solution]
\begin{enumerate}
\item Sans perte de generalite, on suppose que $g\neq 0$. Grace a la division euclidienne on peut ecrire $f=qg+r$ avec $q,r\in K[x]$ et $\deg r< \deg g$, on remarqiue ainsi que $D_{f,g}=D_{g,r}$, pourquoi? En effet soit $h\in K[x]$ alors $h \mid f$ et $h \mid g$ puis $h \mid g$ et $h \mid r$. En continuant ainsi (par des divisions euclidiennes succesives), on arrive a l'existence de $r^{*}\in K[x]$ tel que $D_{f,g}=D_{r^{*},0}$. Cette procedure se termine car la somme de degres de polynomes est strictement decroissante. On a
$$
D_{f,g}=D_{r^{*},0}=\{ d\in K[x]:d\mid r^{*} \}
$$
et le degre des polynomes dans $D_{f,g}$ est borne par $\deg r^{*}$. Soit $a_{n}\in K$ le coefficient dominant de $r^{*}$, on pose
$$
d^{*}:= \frac{1}{a_{n}}r^{*}
$$
et on observe que $d^{*}$ est unitaire et de degre maximal par construction.

Pour monter l'unicite, on prend $\tilde{d}\in D_{f,g}$ de degre maximal et unitaire alors $\tilde{d}\in D_{r^{*},0}=D_{d^{*},0}$ qui implique qu'il existe $\lambda \in K$ tel que $\lambda \tilde{d}=d^{*}$ (car $\deg \tilde{d} = \deg d^{*}$). En comparant les coefficients dominants de $\lambda \tilde{d}$ et $d^{*}$ on conclut que $\lambda=1$ et donc $\tilde{d}=d^{*}$ qui demontre l'unicite de $d^{*}$.
\item Il suffit de montrer que tout diviseur commun de $f$ et $g$ divise $d^{*}$, car $D_{f,g}=D_{d^{*},0}$, et que $d^{*}$ divise $f$ et $g$ car $d^{*}\in D_{f,g}$. On conclut par l'unicite de $d^{*}$.
\end{enumerate}
\end{proof}

\chapter{Serie 4}

\begin{exercise}[Exo plus]
Soit $V$ un espace vectoriel sur un corps $K$ de dimension finie et $f:V\to V$ un endomorphisme. Soit $p:V\to V$ un automorphisme (c-a-d un endomorphisme inversible). Montrer que $\lambda \in K$ est une valeur propre de $f$ si et seulement si $\lambda$ est une valeur propre de l'endomorphisme $p ^{-1}\circ f \circ p$.
\end{exercise}
\begin{proof}[Solution]
\begin{itemize}
\item $\Rightarrow$: Supposons que $\lambda \in K$ est une valeur propre de $f$ avec $\boldsymbol{v}\in V\setminus \{ \boldsymbol{0} \}$ un vectuer propre associe. Posons $\boldsymbol{w}:=p ^{-1}(\boldsymbol{v})\neq \boldsymbol{0}$ (non nul car $p$ est un automorphisme). Si on applique la composition
\begin{align*}
(p ^{-1} \circ f \circ p)(\boldsymbol{w}) & =(p ^{-1} \circ f \circ p)(p ^{-1}(\boldsymbol{v})) \\
 & =p ^{-1} (f(\boldsymbol{v})) \\
 & = p ^{-1} (\lambda \boldsymbol{v}) \\
 & =\lambda p ^{-1} (\boldsymbol{v}) \\
 & =\lambda \boldsymbol{w}
\end{align*}

\item $\Leftarrow$: Supposons que $\lambda \in K$ est un valeur propre de $p ^{-1} \circ f \circ p$ et $\boldsymbol{w}\in V\setminus \{ \boldsymbol{0} \}$ est un vecteur propre associe. Posons $\boldsymbol{v}=p(\boldsymbol{w})\neq 0$ donc
\begin{align*}
f(\boldsymbol{v}) & =f(p(\boldsymbol{w})) \\
 & =(p \circ p ^{-1})(f(p(\boldsymbol{w}))) \\
 & =p(p ^{-1} \circ f \circ p)(\boldsymbol{w}) \\
 & = p(\lambda \boldsymbol{w}) \\  & =\lambda p(\boldsymbol{w})\\
 & = \lambda \boldsymbol{v}
\end{align*}
\end{itemize}
\end{proof}

\begin{exercise}[Exo etoile]
Soient $K$ un corps, et $\lambda_{1},\dots,\lambda _{r}\in K$ les valeurs propres d'une matrice $A\in K^{n\times n}$ et $m_{1},\dots,m_{r}$, leurs multiplicites algebriques. Soit
\begin{equation}
m_{1}+\dots +m_{r}=n \label{hyp51}
\end{equation}

Rappelez-vous que la trace de $A$ est definie par $\mathrm{Tr}(A):=\sum_{i=1}^{n}a_{ii}$.

Demontrer les assertions suivantes:
\begin{enumerate}
\item $\det(A)=\prod_{i=1}^{r}\lambda _{i}^{m_{i}}$
\item si $p_{A}(\lambda)=\alpha _{n}\lambda^{n}+\alpha _{n-1}\lambda^{n-1}+\dots +\alpha_{1}\lambda +\alpha_{0}$ alors $\alpha _{n-1}=(-1)^{n-1}\mathrm{Tr}(A)$
\item $\mathrm{Tr}(A)=\sum_{i=1}^{r}m_{i}\lambda _{i}$
\end{enumerate}
\end{exercise}
\begin{proof}[Solution]
L'hypothese \ref{hyp51} implique que la factorisation en irreductibles du polynome caracteristique est donne par $(-1)^{n}\prod_{i=1}^{r}(x-\lambda _{i})^{m_{i}}=p_{A}(X)=\det(A-xI)$.  
\begin{enumerate}
\item 
\begin{align}
\det(A) & =\det(A-0I) \\
 & =p_{A}(0) \\
 & =(-1)^{n}\prod_{i=1}^{r} (0-\lambda _{i})^{m_{i}}=\prod_{i=1}^{r} \lambda _{i}^{m_{i}}
\end{align}
\item On se rappelle de la formmule de Leibniz (pour les determinants) :
$$
\det(A-xI)=\sum _{\sigma \in S_{n}}\mathrm{sgn}(\sigma)\prod_{j=1}^{n}(A-xI)_{j\sigma (j)} 
$$
comme on s'interese au coefficient $\alpha _{n-1}$, la seule permutation qui est interessante est l'identite. Si $\sigma\neq \mathrm{Id}$, on peut trouver $j\in \{ 1,\dots ,n\}$ tel que $\sigma(j)\neq j$. Pour ce $j$ on a que $(A-xI)_{j\sigma(j)}=A_{j\sigma(j)}$. En plus $\sigma(j)\neq\sigma^{2}(j)$ car $\sigma$ est bijective donc on a $(A-xI)_{\sigma(j)\sigma^{2}(j)}=A_{\sigma(j)\sigma^{2}(j)}$. Prenons donc $\sigma=\mathrm{Id}$ et considerons $$\mathrm{sgn}(\sigma)\prod_{j=1}^{n}(A-xI)_{jj}=\prod_{j=1}^{n} (A_{jj}-x)$$, le coefficient devant $x^{n-1}$ dans ce terme vaut $(-1)^{n-1}\sum_{j=1}^{n}A_{jj}=(-1)^{n-1}\mathrm{Tr}(A)$.
\item On compare les coefficients $\alpha _{n-1}$ obtenu en 2 avec le coefficient de $x^{n-1}$ dans 
$$
(-1)^{n}\prod_{i=1}^{r}(x-\lambda _{i})^{m_{i}} 
$$
Le coefficient de $x^{n-1}$ dans cette formule est donne par
$$
(-1)^{n}\sum_{i=1}^{r} m_{i}(-\lambda _{i})=(-1)^{n+1}\sum_{i=1}^{r} m_{i}\lambda _{i}
$$
en comparant 
$$
(-1)^{n-1}\mathrm{Tr}(A)=(-1)^{n+1}\sum_{i=1}^{r} m_{i}\lambda _{i}\Leftrightarrow \mathrm{Tr}(A)=\sum_{i=1}^{r} m_{i}\lambda _{i}
$$
\end{enumerate}
\end{proof}

\begin{exercise}[Examen 2022]
Quel est le polynome caracteristique de $A\in \mathbf{Q}^{5\times5}$ avec 
$$
A=\begin{pmatrix}
3 & 1 & 2 & 2 & 0 \\
1 & 3 & 0 & 0 & 0 \\
1 & 1 & 1 & 1 & 1 \\
1 & 1 & 2 & 0 & 0 \\
2 & 2 & 2 & 2 & 1
\end{pmatrix}
$$
\end{exercise}
\begin{proof}[Solution]
Regarder coefficients, deviner une valeur propre !
\end{proof}
\chapter{Serie 5}

\begin{exercise}[Exo plus]
Soit $V$ un espace vectoriel et soit $B=\{ v_{1},v_{2},v_{3} \}$ une base de $V$. Soit $f:V\times V\to K$ une forme bilineaire symetrique telle que 
$$
f(v_{1},v_{1})=2,~f(v_{2},v_{2})=3,~f(v_{3},v_{3})=-1,~f(v_{1},v_{2})=0,~f(v_{2},v_{3})=1,~f(v_{3},v_{1})=-2
$$
\begin{enumerate}
\item Ecrire la matrice $A_{B}^{f}$ de la forme bilineaire $f$ pour la base $B$.
\item Trouver une base orthogonale $C=\{ w_{1},w_{2},w_{3} \}$ pour $V$ par rapport a la forme bilineaire $f$.
\end{enumerate}
\end{exercise}
\begin{exercise}[Exo etoile]
Soit $K$ un corps de caracteristique differente de 2, et $V$ un $K$-espace vectoriel de dimension finie $n$. Considerons une forme bilineaire $f$ sur $V$ qui est \textbf{antisymetrique}, c'est a dire 
$$
f(x,y)=-f(y,x), \quad \forall x,y\in V
$$

Montrer que si $f$ est non-degeneree, alors $n$ est pair.

L'implication inverse est-elle vraie?
\end{exercise}
\begin{proof}[Solution]
Fixon $B=\{ b_{1},\dots ,b_{n} \}$ une base de $V$ et considerons la matrice $A_{B}^{f}\in K^{n\times n}$ de la forme bilineaire $f$ dans $B$. Comme $f$ est antisymetrique, on remarque que $A_{B}^{f}$ est aussi antisymetrique c-a-d $(A_{B}^{f})^{\mathsf{T}}=-A_{B}^{f}$. En plus comme $f$ est non-degeneree, par le cours on sait que $\mathrm{rang}(A_{B}^{f})=n$ donc $\det(A_{B}^{f})\neq 0$. On calculer
\begin{align*}
\det(A_{B}^{f}) & =\det((A_{B}^{f})^{\mathsf{T}}) \\
 & = \det(-A_{B}^{f}) \\
 & =(-1)^{n}\det(A_{B}^{f})
\end{align*}
donc comme $\det(A_{B}^{f})\neq0$ et $\mathrm{char}(K)\neq 2$ on a que $(-1)^{n}=1$ donc $n$ est pair.

Si $n$ est pair et $f$ est antisymetrique, alors $f$ n'est pas forcement non-degeneree. Comme contre-exemple on prend la forme nulle $f:V\times V\to K, (v,w)\mapsto 0$.
\end{proof}


\chapter{Serie 6}

\begin{exercise}[Exo plus]
Transformez les matrices reeles suivantes en matrices diagonales dont les elements sont 0, 1 et $-1$. Combien de 0, 1 et $-1$ sont sur la diagonale ? (Ces quantites sont appellees l'indice de nullite, l'indice de positivite et l'indice de negativite.)
$$
\begin{pmatrix}
1 & 2 \\
2 & -1 
\end{pmatrix},~~\begin{pmatrix}
1 & 1 \\
1 & 1
\end{pmatrix},~~\begin{pmatrix}
2 & 4 & 2 \\
4 & 3 & 1 \\
2 & 1 & 1
\end{pmatrix}
$$
\end{exercise}
\begin{proof}[Solution]
\begin{itemize}
\item Pour la premiere matrice on additione la premiere ligene fois $-2$ a la deuxieme ligne pour obtenir
$$
\begin{pmatrix}
1 & 2 \\
0 & -5
\end{pmatrix}
$$
puis on fait le meme pour les colonnes 
$$
\begin{pmatrix}
1 & 0 \\
0 & -5
\end{pmatrix}
$$
et on divise la deuxieme ligne et la deuxieme colonne par $5$
$$
\begin{pmatrix}
 1 & 0 \\
0 & -1
\end{pmatrix}
$$
donc l'indice de positivite est 1, l'indice de negativite est 1 et l'indice de nullite est 0.
\item On soustrait a la deuxieme ligne une fois la premiere
$$
\begin{pmatrix}
1 & 1 \\
0 & 0
\end{pmatrix}
$$
et de meme pour les colonnes
$$
\begin{pmatrix}
1 & 0 \\
0 & 0
\end{pmatrix}
$$
donc on obtient l'indice de positivite 1, l'indice de negativite 0 et l'indice de nullite 1.
\item On effectue des operations elementaires sur les lignes et colonnes 
$$
\begin{pmatrix}
2 & 4 & 2 \\
4 & 3 & 1 \\
2 & 1 & 1
\end{pmatrix} \underset{L_{2}-2L_{1}}\longrightarrow \begin{pmatrix}
2 & 4 & 2 \\
0 & -5 & -3 \\
2 & 1 & 1
\end{pmatrix} \underset{C_{2}-2C_{1}}\longrightarrow \begin{pmatrix}
2 & 0 & 2 \\
0 & -5 & -3 \\
2 & -3 & 1
\end{pmatrix}
$$
puis
$$
\underset{L_{3}-L_{1}}\longrightarrow \begin{pmatrix}
2 & 0 & 2 \\
0 & -5 & -3 \\
0 & -3 & -1
\end{pmatrix} \underset{C_{3}-C_{1}}\longrightarrow \begin{pmatrix}
2 & 0 & 0 \\
0 & -5 & -3 \\
0 & -3 & -1
\end{pmatrix}
$$
puis
$$
\underset{L_{3}-\frac{3}{5}L_{2}}\longrightarrow \begin{pmatrix}
2 & 0 & 0 \\
0 & -5 & -3 \\
0 & 0 & \frac{4}{5}
\end{pmatrix} \underset{C_{3}-\frac{3}{5}C_{2}}\longrightarrow \begin{pmatrix}
2 & 0 & 0 \\
0 & -5 & 0 \\
0 & 0 & \frac{4}{5}
\end{pmatrix}
$$
puis on divise les lignes par 2 , 5 et $\frac{5}{4}$ respectivement pour obtenir l'indice de positivite 2, l'indice de negativite 1 et l'indice de nullite 0.
\end{itemize}
\end{proof}
\begin{exercise}[exo 7]\label{exo7}
Soit $g:V\to V$ diagonalisable et $W\subseteq V$ est invariant par $g$. Alors $g|_{W}:W\to W$ est diagonalisable.
\end{exercise}
\begin{proof}[Solution]
Comme $g$ est diagonalisable on a $V=\bigoplus _{\lambda}E_{\lambda}$. Pour montrer que $g|_{W}$ est diagonalisable aussi, il suffit d'observer que $\bigoplus _{\lambda}(W\cap E_{\lambda})=W$, on remarque facilement que la direction / inclusion $\subseteq$ tient. 

Montrons $\supseteq$, prenons $w\in W$ et observons que $w\in \bigoplus_{\lambda}(W\cap E_{\lambda})$. On ecrit $w=\sum _{i} v_{\lambda _{i}}$ ou $v_{\lambda _{i}}\in E_{\lambda _{i}}$ et pour $i\neq j$ on a $\lambda _{i}\neq\lambda _{j}$. Pour tout $i$ on considere 
\begin{align*}
\underbrace{ \prod _{j\neq i} \left( \frac{g-\lambda _{j}\mathrm{Id}}{\lambda _{i}-\lambda _{j}} \right) }_{ h }(w) & =\prod _{j\neq i} \left( \frac{g-\lambda _{j}\mathrm{Id}}{\lambda _{i}-\lambda _{j}} \right)\left( \sum _{k}v_{\lambda _{k}} \right) \\
 & =v_{\lambda _{i}}\in W
\end{align*}
car $h(v_{\lambda _{j}})=0$ pour $j\neq i$ et $h(v_{\lambda _{j}})=v_{\lambda _{i}}$ pour $j=i$. Donc $w\in \bigoplus _{\lambda}(W\cap E_{\lambda})$.
\end{proof}

\begin{exercise}[Exo etoile]
Monter que si $f$ et $g$ sont diagonalisables et commutent, alors $f$ et $g$ sont simultanement diagonalisables.
\end{exercise}
\begin{proof}[Solution]
Comme $f$ est diagonalisable, on peut ecrire
\begin{equation}
V=\bigoplus_{\lambda} E_{\lambda} \label{basestar}
\end{equation}
ou les $\lambda$ sont les valeurs propres de $f$ et $E_{\lambda}$ les espaces propres de $f$ (c-a-d $E_{\lambda}=\ker(f-\lambda \mathrm{Id})$). Montrons que $E_{\lambda}$ est $g$-invariant. Prenons $v\in E_{\lambda}$ et montrons que $g(v)\in E_{\lambda}$:
\begin{align*}
(f-\lambda \mathrm{Id})(g(v)) & =(f\circ g-\lambda g )(v) \\
 & =(g\circ f- \lambda g)(v) \\
 & =g(f-\lambda \mathrm{Id} )(v) \\
 & =g(0)=0
\end{align*}
donc $E_{\lambda}$ est $g$-invariant pour tout $\lambda$ une valeur propre de $f$.

Du coup par l'exercise 7 de la serie 6 \ref{exo7} pour chaque valeur propre $\lambda$ de $f$ on trouve une base de $B_{\lambda}$ de $E_{\lambda}$ forme de vecteurs propres par rapport a $g$ (car $g$ est diagonalisable, $E_{\lambda}$ est $g$-invariant donc $g|_{E_{\lambda}}$ est diagonalisable). Finalement, 
$$
B=\bigcup _{\lambda}B_{\lambda}
$$
est une base de $V$ (a cause de \ref{basestar} et du fait que $B_{\lambda}$ est une base de $E_{\lambda}$), en plus $B$ consiste de vecteurs propres a la fois de $f$ et de $g$ (les elements de $B$ vivent dans les espaces propres de $f$ et sont par construction des vecteurs propres pour $g$). 
\end{proof}
\chapter{Serie 7}
\begin{exercise}[Exo plus]
Soit $A\in \mathbf{C}^{n\times n}$ une matrice hermitienne et soit $\lambda$ une valeur popre de $A$. Montrer que $\lambda$ est reele.
\end{exercise}
\begin{proof}[Solution]
Sans preuve.
\end{proof}

\begin{exercise}[Exo etoile]
Le but de cet exercice est de montrer \emph{l'inegalite d'Hadamard}, c'est-a-dire
$$
|\det(A)|\le \prod_{i=1}^{n} \lVert a_{i} \rVert_{2},
$$
pour $A\in \mathbf{R}^{n\times n}$ dont les colonnes sont $a_{1},\dots,a_{n}$.

Soit $A'$ la matrice dont les colonnes sont $a_{1}^{*},\dots,a_{n}^{*}$, les vecteurs \emph{non-normalises} issus de l'orthogonalisation de Gram-Schmidt sur les colonnes de $A$. On rappelle la décomposition $A=A'S$, ou $S$ est triangulaire supérieure de diagonale 1.
\begin{enumerate}
\item Montrer que $\det(S)=1$ et que $\det(A)=\det(A')=\pm \prod_{i=1}^{n}\lVert a_{i}^{*} \rVert_{2}$.
\item Montrer que $\lVert a_{i}^{*} \rVert\le \lVert a_{i} \rVert$ pour tout $i=1,\dots,n$ et conclure.
\end{enumerate}

Supposons de surcroit que les coefficients de $A$ soient tous bornés absolument par $M\in \mathbf{R}_{+}:|A_{ij}|\le M$ pour tous $i$, $j$. Déduire de l'inégalité d'Hadamard que $|\det(A)|\le M^{n}n^{n/2}$.
\end{exercise}
\begin{proof}[Solution]
\begin{enumerate}
\item Comme $S$ est triangulaire supérieure avec des coefficients diagonaux égaux à 1, on obtient $\det(S)=1$. On s'intéresse à la matrice $A'^{\mathsf{T}}A'$ :
\begin{align*}
A'^{\mathsf{T}}A' & =\begin{pmatrix}
(a_{1}^{*})^{\mathsf{T}}a_{1}^{*} & (a_{1}^{*})^{\mathsf{T}}a_{2}^{*} & \cdots & (a_{1}^{*})^{\mathsf{T}}a_{n}^{*} \\
(a_{2}^{*})^{\mathsf{T}}a_{1}^{*} & \cdots & \cdots & \vdots \\
\vdots &  &  & \vdots \\
(a_{n}^{*})^{\mathsf{T}}a_{n}^{*} & \cdots & \cdots & (a_{n}^{*})^{\mathsf{T}}a_{n}^{*}
\end{pmatrix} \\
 & =\mathrm{diag}(\lVert a_{i}^{*} \rVert_{2}^{2},\dots,\lVert a_{n}^{*} \rVert _{2}^{2} )\footnotemark{}
\end{align*}
\footnotetext{Justifie par la procedee de Gram-Schmidt}
Donc
\begin{align*}
\det(A')^{2} & =\det(A')\det(A') \\
 & =\det(A'^{\mathsf{T}})\det(A') \\
 & =\det(A'^{\mathsf{T}}A') \\
 & =\prod_{i=1}^{n} \lVert a_{i}^{*} \rVert _{2}^{2}
\end{align*}

Comme $\det(A)=\det(A')\det(S)=\det(A')\cdot 1$ alors $$\det(A)=\det(A')=\pm \prod_{i=1}^{n}\lVert a_{i}^{*} \rVert_{2}.$$
\item On a
\begin{align*}
\lVert a_{i} \rVert _{2}^{2} & =\lVert a_{i}-a_{i}^{*} +a_{i}^{*}  \rVert_{2}^{2} \\
 & =\lVert a_{i}-a_{i}^{*} \rVert _{2}^{2} +\lVert a_{i}^{*} \rVert _{2}^{2} \\
 & \ge \lVert a_{i}^{*} \rVert _{2}^{2} 
\end{align*}

On justifie le fait que $a_{i}^{*}\perp a_{i}-a_{i}^{*}$. Par Gram-Schmidt, on a
\begin{align*}
a_{i}^{*} & =a_{i}-\frac{\langle a_{i},a_{1}^{*} \rangle }{a_{1}^{*},a_{1}^{*}}a_{1}^{*}-\dots - \frac{\langle a_{i},a_{i-1}^{*} \rangle }{\langle a_{i-1}^{*},a_{i-1}^{*} \rangle }a_{i-1}^{*}
\end{align*}
donc $a_{i}-a_{i}^{*}\in \mathrm{span}\{ a_{1}^{*},\dots,a_{i-1}^{*} \}$ qui est perpendiculaire à $a_{i}^{*}\in \mathrm{span}\{ a_{1}^{*},\dots,a_{i-1}^{*} \}$.
\end{enumerate}

Si les coefficients sont bornées par $M\in \mathbf{R}_{+}$, on a 
\begin{align*}
\lVert a_{i} \rVert _{2} & =\sqrt{ (a_{i})_{1}^{2}+\dots +(a_{i})_{n}^{2} } \\
 & \le \sqrt{ M^{2}+\dots +M^{2} } \\
 & =\sqrt{ n }M
\end{align*}

Par l'inégalité d'Hadamard, on a 
\begin{align*}
|\det(A)| & \le \prod_{i=1}^{n} \lVert a_{i}\rVert_{2} \\
 & \le \prod_{i=1}^{n} \sqrt{ n }M \\
  & =M^{n}n^{n/2}.
\end{align*}
\end{proof}

\chapter{Serie 8}

\begin{exercise}[Exo plus]
Soit $A\in \mathbf{C}^{n\times n}$ une matrice hermitienne et $u,v\in \mathbf{C}^{n}$ deux vecteurs propres associés a des valeurs propres différentes. Montrer que $u$ et $v$ sont orthogonaux par rapport au produit hermitien standard.
\end{exercise}
\begin{proof}[Solution]
Soient $\lambda_{1}\neq\lambda_{2}$ deux valeurs propres distinctes et $u,v\in \mathbf{C}^{n}\setminus \{ 0 \}$ leurs vecteurs propres associés respectivement. Si $u^{*}v\neq 0$, alors $\lambda_{1}u^{*}v\neq\lambda u^{*}v$. On a
\begin{align*}
\lambda_{1}u^{*}v & =(\lambda_{1}u)^{*}v \\
 & =(Au)^{*}v \\
 & =u^{*}A^{*}v \\
 & =u^{*}Av \\
 & =\lambda_{2}u^{*}v
\end{align*}
alors $u^{*}v=0$.
\end{proof}
\begin{remark}
Une autre definition de matrice hermitienne : soit $A\in \mathbf{C}^{n\times n}$ alors $A^{\mathsf{T}}=\overline{A}$.
\end{remark}
\end{document}